\section{Conclusion}
\label{sec:conclusion}

This study finds that Indonesian \id{-nya} preprocessing has small retrieval effects on MIRACL-id. BM25 is statistically sensitive to strategy choice, but practical gaps are tiny; BGE-m3 is nearly invariant. The safest default is to preserve raw text for dense retrieval and use dictionary-guarded clitic handling, not naive regex stripping, when sparse pipelines require suffix handling. The in-house resolver does not improve retrieval, leaving stronger Indonesian coreference models and multi-benchmark replication as future work. The released code, metric CSVs, analysis outputs, and integrity reports make this negative result reproducible.

\section*{Reproducibility}

All code, processed data, intermediate indices, and analysis scripts are available at \url{https://github.com/shaneemichael/TA_TBI}. The methodology was pre-registered via a timestamped commit before experimental code was run; the relevant commit hash is reported in the supplementary materials.

\section*{AI assistance disclosure}

Literature scanning, code skeleton drafting, and LaTeX scaffolding were assisted by Claude (Sonnet/Opus). All experimental design choices, hypothesis formulation, statistical analyses, interpretation of results, and final manuscript content are the work of the human author.

\section*{Acknowledgements}

We thank the MIRACL dataset creators for releasing the Indonesian retrieval benchmark, the Sastrawi maintainers for the Indonesian stemming resources used in the dictionary-guarded conditions, and the open-source maintainers of Pyserini, FAISS, FlagEmbedding, and pytrec-eval.
